% !TeX program = pdflatex
% ALIUS Bulletin standalone fictional parody source.
% Compile from this directory with pdflatex for the final PDF.
\documentclass[a4paper,11pt]{article}
\usepackage[margin=22mm]{geometry}
\usepackage[T1]{fontenc}
\usepackage[utf8]{inputenc}
\usepackage{xcolor}
\usepackage[hidelinks]{hyperref}
\usepackage[protrusion=true,expansion=false]{microtype}
\usepackage{parskip}
\usepackage{fancyhdr}

\definecolor{ALIUSRed}{RGB}{168,28,38}
\definecolor{ALIUSGrey}{RGB}{75,75,75}

\pagestyle{fancy}
\fancyhf{}
\fancyhead[L]{\footnotesize ALIUS Bulletin fictional supplement}
\fancyhead[R]{\footnotesize Peterson - Thermodynamic Superego of the Lobster}
\fancyfoot[C]{\footnotesize\thepage}
\renewcommand{\headrulewidth}{0.4pt}
\setlength{\headheight}{14pt}

\newenvironment{question}{\par\noindent\color{ALIUSRed}\textbf{Zizek:}\ }{\par}
\newenvironment{answer}{\par\noindent\color{black}\textbf{Peterson:}\ }{\par}

\title{\textbf{The Thermodynamic Superego of the Lobster}\\
\large A fictional interview on consciousness, hierarchy, entropy, and pure ideology}
\author{A fictional interview with Jordan B. Peterson by Slavoj Zizek}
\date{ALIUS Bulletin fictional supplement, 2026}

\begin{document}
\maketitle



\section*{Metadata}
\textbf{Piece type:} fictional interview / parody\par
\textbf{Citation:} Peterson, J. B., and Zizek, S. (fictional). (2026). The thermodynamic superego of the lobster: A fictional interview on consciousness, hierarchy, entropy, and pure ideology. \textit{ALIUS Bulletin fictional supplement}.\par
\textbf{Keywords:} consciousness; altered states; hierarchy; entropy; psychoanalysis; postmodernism; ideology; lobsters\par

\section*{Abstract}
In this fictional interview, Slavoj Zizek and Jordan B. Peterson are staged inside an imaginary ALIUS dialogue on consciousness, hierarchy, entropy, postmodernism, and ideology. Zizek's questions proceed by way of Hegel, Lacan, Marx, cinema, class consciousness, and the psychoanalytic status of the lobster, without ever arriving at a stable point of accusation. Peterson's answers respond with maximal confidence, anatomical solemnity, mythological escalation, and a dense vocabulary of antientropic order, archetypal recurrence, serotonergic domination, and symbolic housekeeping. The result is a deliberately inflated conversation in which the surface resembles scholarship, the affect resembles intellectual combat, and the underlying content behaves like a fog machine that has read too much German idealism. The interview asks, in its own absurd mode, whether consciousness is a ladder, a mistake, a political symptom, a thermodynamic rebellion, or simply what happens when a nervous system becomes sufficiently worried about its room.

\section*{Interview}

\begin{question}
Let us begin, if you agree, with the lobster, but not, as people too quickly imagine, with the lobster as this stupid biological example, the little marine monarch waving its claws in the restaurant aquarium. I am interested in the lobster as what Lacan might have called the obscene crustacean supplement of the symbolic order: the animal that appears when the big Other wants hierarchy but does not want to admit that hierarchy is already an ideology of the body. So, my question is almost embarrassingly simple, and therefore, of course, completely impossible. When you speak of lobster hierarchies as evidence that dominance precedes capitalism, postmodern social construction, and perhaps even the liberal dream of infinite plasticity, are you describing consciousness as a biological ladder, or are you describing the ladder itself as the fantasy through which consciousness misrecognizes its own fall?
\end{question}

\begin{answer}
Well, it is a very serious question, and it is serious precisely because it begins in the wrong place, which is often where the most serious questions begin. You see, the lobster is not merely a lobster. That is the first thing. If you think the lobster is merely a lobster, then you have already collapsed the ontological hierarchy of being into a flattened, resentful, quasi-Marxist taxonomy of edible objects, and that is a very dangerous thing to do, especially if you are young and your room is in a state of symbolic and literal disarray.

The lobster is a phylogenetically ancient signal, a kind of walking neurological footnote, or perhaps I should say swimming neurological footnote, to the brutal and beneficent fact that being organizes itself vertically before it organizes itself ethically. That does not mean the vertical is always good. People make this mistake. They say, "Peterson says hierarchy exists, therefore Peterson worships hierarchy." No. That is absurd. That is like saying because I notice gravity I am in favor of falling down the stairs. But if you pretend gravity is merely a discourse, you will still meet the stairs, and the stairs will not have read Foucault.

Now, consciousness enters this picture as the dramatic theater in which an organism encounters the unbearable necessity of rank, aim, and sacrifice. The nervous system does not wake up in an egalitarian void. It wakes up in a field of affordances, threats, ancestral memories, maternal expectations, paternal prohibitions, and the terrible fact that the future is coming at you whether you have made yourself intelligible to it or not. That is what the lobster knows, in its inarticulate but neurochemically profound way. It knows that defeat curls the body inward and victory lifts the body upward. That is not an ideology. That is posture as metaphysics.

But to your deeper point, the ladder is indeed a fantasy, and this is where people think I am contradicting myself, but I am not. The ladder is a fantasy because every representation of order is symbolic. The moment you say "up," you have already imported a mythology of ascent, a cathedral architecture, Jacob's ladder, the axis mundi, the terrible mountain, the dominance hierarchy, the child's report card, the university tenure committee, and your grandmother asking why you are not married. All of that is in the word "up." But it is not only fantasy. It is the necessary fantasy by which finite beings orient themselves in an entropic universe.

So consciousness is not simply a ladder. It is the capacity to suffer the ladder, to resent it, to climb it, to discover that it was partly imaginary, and then to build a better one while apologizing to everyone you stepped on because you confused ambition with revelation. That is a very difficult thing. It is much more difficult than saying the ladder is oppressive and then lying on the floor in an ethical tantrum. The lobster does not solve that. But it reminds us, with a kind of prehistoric comedic solemnity, that before ideology speaks, the body has already taken a position.
\end{answer}

\begin{question}
Yes, yes, but here we should be careful. I would even say, if I may be vulgar in a properly philosophical way, that the lobster is not simply before ideology, because this "before" is already the retroactive fantasy of ideology itself. The commodity also says: I am just a thing, please buy me. And then, of course, precisely in this humility, it hides the whole universe of exploitation. So maybe the lobster says: I am just serotonin, I am just hierarchy, no politics here. But perhaps this "just" is where the symbolic smuggles itself in. How do you distinguish the biological substrate of rank from the ideological enjoyment we derive from pretending the substrate has excused us from moral responsibility?
\end{question}

\begin{answer}
You distinguish them badly at first, and then, if you are fortunate and disciplined and not entirely possessed by your own cleverness, you distinguish them a little better. That is the honest answer. The human animal is not capable of a clean separation between substrate and symbol. We are embodied metaphors. The body throws up patterns from deep time, and consciousness clothes those patterns in stories, and then the stories start giving orders to the body. That circularity is not a bug in the system. It is the system.

Now, the question of moral responsibility is absolutely central. Biology does not excuse tyranny. Biology explains the arena in which tyranny tempts us. The fact that hierarchies are ancient does not mean that every hierarchy is justified. The fact that dominance has a serotonin-mediated component does not mean that whoever has the biggest office has achieved spiritual enlightenment. That would be foolish. It would also be bureaucratically convenient, which is one reason institutions enjoy pretending it is true.

The danger on the ideological left, if I can put it that way, is the tendency to see all rank as domination and all domination as arbitrary. That is a catastrophic oversimplification. It blinds you to competence, and competence is the civilized form of dominance. You want your surgeon to be high in a hierarchy of surgical competence. You want your airplane pilot to be high in a hierarchy of not confusing the runway with a poststructuralist invitation to reinterpret descent. There are hierarchies that are approximations of value, and if you destroy them all in the name of an undifferentiated compassion, you do not get equality. You get chaos, and then you get a worse hierarchy.

But the danger on the right, which is equally real, is to sacralize the hierarchy because it is old, as if age were the same thing as wisdom. Cancer is old. Parasites are old. The capacity of a chimpanzee to rip another chimpanzee apart is old. Old is not enough. What matters is whether the hierarchy has been transformed by truth, reciprocity, responsibility, and the tragic recognition that every ascent must be answerable to what it leaves below.

So when you say ideological enjoyment, I would translate that into the sinful pleasure of reducing moral complexity to a permission slip. The resentful egalitarian says, "All hierarchy is oppression, therefore I owe nothing to excellence." The authoritarian says, "Hierarchy is natural, therefore I owe nothing to compassion." Both positions are infantile. Both avoid the burden of consciousness. Consciousness is the terrible mediator that says: yes, you are an animal; yes, you are symbolic; yes, you inherit dominance; yes, you must answer for what you do with it. The lobster cannot answer for its hierarchy. We can. That is why our claws are verbal, and that is why our sins have footnotes.
\end{answer}

\begin{question}
Now, let us risk the big word: consciousness. In some theories, consciousness is a global workspace, in others a predictive inference process, in others a strange loop, and in yet others the minimum necessary illusion a nervous system produces so that it can blame itself in the morning. But if we follow Hegel, consciousness is never simply there. It is what discovers, too late, that its certainty was already mediated. And if we follow Lacan, the subject is not the captain of the ship but the leak in the ship, the hole around which the ship is organized. So, when you speak of consciousness in relation to responsibility, are you not making consciousness too heroic? What if consciousness is not the hero facing chaos, but the little bureaucrat inside chaos stamping documents after the catastrophe has already happened?
\end{question}

\begin{answer}
That is a very clever formulation, and it contains a great deal of truth, although perhaps it hides from the truth it contains by wrapping itself in theatrical modesty. Consciousness is not always heroic. In fact, most of the time consciousness is pathetic, evasive, post hoc, self-serving, and late. Anyone who has watched themselves make the same mistake for the four hundredth time understands that. We often do not decide so much as narrate the inevitability of our failure afterward in a tone that suggests metaphysical sophistication.

But heroism is not the absence of belatedness. Heroism is the decision to enter the catastrophe retroactively and behave as if the future could still be redeemed. That is the paradox. The hero is always late. You do not need a hero before the dragon arrives. You need one when the village is already burning and everyone has discovered that the dragon was partly external, partly familial, partly economic, partly libidinal, and partly the consequence of not doing the laundry for seven months.

The models you mention are useful, each in their own way. The global workspace gives us a picture of consciousness as broadcast, a theater of integration. Predictive processing tells us that perception is not passive reception but active anticipation constrained by error. The strange loop captures the reflexive, self-referential structure of the subject, the snake eating its tail while publishing a monograph about digestion. But these models, though powerful, risk becoming bloodless if they do not confront the moral texture of conscious life. Consciousness is not merely information made available to systems. It is information made answerable to suffering.

That is why responsibility matters. Not because consciousness is a sovereign king in the skull, issuing decrees from a Cartesian balcony. That is naive. Responsibility matters because consciousness is the trembling point at which automaticity can be interrupted, even slightly. You are right that the bureaucrat often arrives late. But sometimes the bureaucrat refuses to stamp the usual form. Sometimes the little clerk in the catastrophe says: no, this time the dragon does not get the permit. That tiny refusal is civilization.

And here is where the heroic language becomes necessary, not as description but as invitation. Human beings do not live by accurate mechanistic models alone. They live by narrative orientation. If you tell people consciousness is merely a post hoc rationalization machine, many will use that as permission to become worse. If you tell them they are potential heroes, some will inflate themselves into absurdity, yes, but others will clean one drawer, tell one truth, apologize to one person, and prevent the future from becoming a repetition compulsion with better branding. So perhaps consciousness is a bureaucrat. But under the right conditions, and with sufficient tragic pressure, even a bureaucrat can become Saint George with a clipboard.
\end{answer}

\begin{question}
This is precisely where I want to complicate, although not to disagree, because disagreement today is too often just narcissism with a microphone. You say: "tell one truth." But truth, for Hegel, is not the little empirical fact standing there like a potato. Truth is the whole, and the whole is false if it arrives as a closed system. And for Marx, truth has a class position even when it thinks it is speaking from nowhere. So, when the individual tells the truth, maybe he is also reproducing the coordinates that make his suffering intelligible only as private failure. Is there not a danger that the ethics of cleaning your room becomes the perfect ideology of late capitalism: everyone privately tidies the symptom while the structure rents out the mess?
\end{question}

\begin{answer}
Of course there is a danger. Everything powerful is dangerous. Water is dangerous if you are drowning. Fire is dangerous if it is in the nursery. Language is dangerous if it is in the hands of someone who has discovered that resentment can be made to sound compassionate. The fact that a principle can be weaponized does not invalidate the principle. It demands discernment.

"Clean your room" is perhaps the most comically misunderstood statement in the modern intellectual landscape, which is saying something because the landscape is littered with comically misunderstood statements. It does not mean: ignore society. It does not mean: capitalism is perfect if your socks are paired. It does not mean: all suffering is your fault and history began when you failed to make your bed. Those are caricatures so cheap they should be taxed as ideological fast food.

The point is developmental and phenomenological. The room is the most immediate domain in which chaos and order meet under your voluntary attention. It is not the whole world, but it is not nothing. It is the first negotiated territory. If you cannot bring order to the sphere where your agency is maximal, you are unlikely to bring order to spheres where your agency is diluted by institutions, law, economics, bureaucracy, and the accumulated stupidity of dead committees. The room is not a substitute for politics. It is the training ground for perception.

Now, your Marxist concern is legitimate. There are structures that produce disorder at a scale no individual can tidy away. Poverty, exploitation, addiction, family collapse, bad schools, corrupt institutions, technological capture, informational chaos, all of these things distribute mess unevenly. A person in a collapsing neighborhood does not face the same room as a person with inherited capital, a quiet street, and parents who read to them at bedtime. Any serious ethic must acknowledge that. But the acknowledgment must not become theft of agency. Compassion that annihilates agency becomes sentimental domination.

The profound problem is that both individual responsibility and structural critique are true, and they are true at different levels of analysis. The immature mind wants one level to devour the other. The ideological mind says: only the individual, or only the structure. But consciousness matures by tolerating multi-level causality without collapsing into paralysis. You clean your room so that you can see the world more clearly. Then, if you are serious, you discover that the hallway is on fire. Then you organize with your neighbors. Then you discover the zoning board is corrupt. Then you discover your own desire for revenge is corrupting your reform. Then you return to your room, not because politics is false, but because the revolutionary who cannot distinguish justice from personal hatred will eventually build a prison and call it health care.
\end{answer}

\begin{question}
Let us move from the room to entropy, because here the mess becomes cosmic. In thermodynamics, entropy is not simply disorder in the vulgar sense, but the statistical tendency toward dispersion, toward the exhaustion of available differences. But ideology also loves entropy. It says: everything becomes fluid, everything becomes exchangeable, everything becomes content, even rebellion. Postmodernism, in its worst caricature, is maybe the heat death of meaning: no privileged signifier, no master narrative, only little temperature differences in the marketplace of identities. Do you think consciousness is an antientropic project? And if yes, is the end of entropy salvation, tyranny, or the most boring possible heaven?
\end{question}

\begin{answer}
The end of entropy would be hell disguised as accounting. Let us be precise. Consciousness is antientropic only in the local and tragic sense. Life itself is a temporary pattern of improbability maintained against dissolution by energy, boundary, metabolism, and sacrifice. A cell is a little cathedral of negentropy. A mind is a cathedral that knows the roof is leaking. Human culture is a cooperative hallucination designed to keep the roof up long enough for children to learn songs.

But you do not want the end of entropy. Entropy is not merely the enemy. Entropy is the condition against which form becomes meaningful. If nothing can decay, nothing can be cared for. If nothing can be lost, courage evaporates. If time cannot wound you, love becomes a scheduling problem among immortals. The antientropic project of consciousness is not to abolish entropy but to negotiate with it without lying. That is why cleaning your room matters at the trivial level and why ritual matters at the civilizational level. Both say: we know chaos is coming, and nevertheless we arrange the flowers.

Now, postmodernism. There is a sophisticated postmodernism and there is the campus-meme postmodernism that wanders around with scissors cutting the wires between competence and recognition. The sophisticated version reminds us that perception is mediated, language is not transparent, power can hide inside categories, and the speaker is not innocent. Fine. That is all useful, and anyone with a decent education should know it. But the pathological version takes mediation to mean arbitrariness, power to mean evil, and interpretation to mean infinite permission. That is where meaning begins to thermally dissolve.

Consciousness requires difference. It requires figure and ground, aim and obstacle, truth and falsehood, self and other, yesterday and tomorrow, lobster and not-lobster. If all differences are reduced to arbitrary constructions, consciousness loses its grip on the world. But if all differences are frozen into eternal categories, consciousness becomes tyrannical. The proper symbolic order is neither liquid nor stone. It is something like a living tree: structured, differentiated, rooted, growing, capable of pruning, and very annoyed if you replace it with a laminated policy document.

So the end of entropy, if imagined as perfect order, is not salvation. It is totalitarian stillness. It is a world in which nothing surprising can appear, no error can teach, no child can interrupt a philosopher, no joke can expose a king. But pure entropy is not freedom either. Pure entropy is noise, dissolution, the inability to bind attention to anything stable enough to love. Consciousness lives between these horrors. It is the sacrificial maintenance of habitable order at the edge of chaos. The edge is not a slogan. It is where being becomes adventure instead of inventory.
\end{answer}

\begin{question}
I like this "edge," although perhaps the edge is also where ideology is most comfortable, because every system today sells itself as dynamic, creative, disruptive, flexible, and so on. Capital says: I am the edge of chaos. The university says: I am criticality. Even the management consultant speaks like Heraclitus with a spreadsheet. So when you say we must live on the edge of chaos, how do we avoid converting myth into motivational management? Is there a difference between the archetype and the TED talk, or is the TED talk the archetype after it has been bitten by the commodity form?
\end{question}

\begin{answer}
That is a beautiful and horrible formulation. The TED talk is the archetype after it has been bitten by the commodity form and then forced to smile under stage lighting for eighteen minutes. I grant you that. But the corruption of an archetype does not erase the archetype. It shows that the archetype is potent enough to be exploited.

When corporations tell you to embrace disruption, they often mean: please metabolize our irresponsibility as your personal growth. That is not the edge of chaos in the mythological sense. That is managerial entropy laundering. The authentic edge of chaos is where your current identity encounters the necessity of transformation in relation to a real good. The fake edge is where an institution externalizes risk downward while retaining the vocabulary of creativity upward. They are not the same. One is Exodus; the other is a quarterly slide deck.

Myth becomes motivational management when it is detached from sacrifice. That is the diagnostic criterion. If the dragon is only an obstacle to your brand, it is not a dragon. If chaos is only a market opportunity, it is not chaos. If transformation costs you nothing except the price of a seminar, you have not transformed; you have purchased symbolic cologne. Real myth is about death and rebirth, and not in the decorative sense. It says: the part of you that is inadequate must actually die, and you will not enjoy that, because the inadequate part of you has friends, habits, playlists, legal justifications, and a very convincing inner lawyer.

The archetype is dangerous because it addresses the preconceptual structures of motivation. It grips you before you have formulated a policy position. That is why totalitarian movements love myth. That is why advertising loves myth. That is why politics loves images of dawn, soil, children, workers, flags, and renovated kitchens. Myth bypasses the committee of propositions and recruits the body. But the answer is not to abandon myth. The answer is to tell truer myths, more self-critical myths, myths that include the shadow of their own misuse.

The heroic archetype, properly understood, is not self-aggrandizement. It is voluntary burden-bearing in the service of what is genuinely good. The TEDified archetype says: become extraordinary so others will applaud your optimized authenticity. The real archetype says: descend into what terrifies you, bring back what the community needs, and do not lie about the monsters you met, especially the ones wearing your face. That distinction matters. Without it, yes, everything degenerates into motivational management, and the dragon becomes a logo with a values statement.
\end{answer}

\begin{question}
Let me introduce psychoanalysis more directly, because perhaps the dragon wearing your face is what Lacan names desire: not what I want, but the way my wanting is structured by the desire of the Other. The subject says: I voluntarily bear the burden. But maybe the burden is already staged by the symbolic network, by the father, by the market, by the professor, by the algorithm that knows I search for "meaning" at 2:13 in the morning. How can responsibility remain responsibility if desire itself is not mine, if even my heroism is caught in the Other's demand?
\end{question}

\begin{answer}
Responsibility begins precisely there. If desire were transparently yours, responsibility would be much simpler and much less necessary. The fact that your desires are inherited, mimetic, distorted, contaminated, wounded, seduced, advertised to, algorithmically amplified, and half-possessed by ghosts does not eliminate responsibility. It deepens it.

You are not responsible for the entire origin of your desire. That would be insane. You did not choose your nervous system, your parents, your historical moment, your temperament, your childhood injuries, your culture's symbolic debris, or the fact that your phone is a slot machine designed by extremely intelligent people who would prefer you not develop a soul. But you are responsible for entering into dialogue with those determinants. That is the narrow gate.

Psychoanalysis, at its best, understands this. It does not say, "You are free in the shallow liberal sense, now choose a breakfast cereal and invent yourself." Nor does it say, "You are determined, therefore collapse elegantly into theory." It says: speak, listen, repeat, notice, suffer, and perhaps the compulsion will become conscious enough to be negotiated. That is a religious process in secular clothing, and I mean that as a compliment.

The heroism I am describing is not pure autonomy. Pure autonomy is a fantasy of adolescents and libertarians with insufficient experience of influenza. Human beings are embedded. They are spoken before they speak. They are loved badly before they love badly. They inherit ancestral scripts, many of which are useful, some of which are demonic, and most of which are ambiguous enough to keep scholars employed indefinitely. Responsibility is the act of discriminating among these scripts under conditions of incomplete knowledge.

So when the Other demands, the first question is not "How do I escape demand?" You cannot. The first question is "Which demand is worthy of my sacrifice?" The father may demand conformity, or discipline, or courage, or silence. The market may demand productivity, spectacle, consumption, or endless self-reinvention. The algorithm may demand arousal without commitment. The body may demand comfort. The conscience may demand truth. These demands conflict. Consciousness is the tribunal in which they are cross-examined, badly at first, and then perhaps with increasing precision.

Your heroism is never uncontaminated. That is why humility is necessary. But contamination does not make heroism fake. It makes it tragic. The purified hero is a propaganda poster. The real hero is partially ridiculous, narcissistic, frightened, resentful, horny, tired, overcaffeinated, historically situated, and still capable of telling the truth. That "still" is the miracle.
\end{answer}

\begin{question}
This "still" is close to what I would call the gap, the crack in determination. But now I must ask about class consciousness, because today altered states are everywhere: psychedelics in therapy, meditation apps, neurofeedback, executive retreats, ecstatic festivals, sleep optimization, microdosing, trauma language, and the corporate appropriation of shamanic vocabulary. Suppose class consciousness itself is an altered state. Not simply knowledge of exploitation, but a transformation in the coordinates of what can be perceived. Can there be a psychedelic Marxism of consciousness, or does the altered state too easily become another bourgeois spa for the suffering ego?
\end{question}

\begin{answer}
There can certainly be altered states that reveal class realities otherwise hidden by ordinary consciousness. There can also be altered states that reveal, with great emotional conviction, absolute nonsense. That is the problem with revelation. Revelation does not come with peer review, and the gods do not always footnote their claims.

Psychedelics, meditation, fasting, ritual, music, grief, love, illness, childbirth, terror, disciplined study, and genuine political solidarity can all alter the structure of salience. They can make previously invisible patterns appear with tremendous force. A worker may suddenly perceive that what he had interpreted as personal inadequacy is partly the consequence of a system designed to extract labor while privatizing blame. That is a real transformation of consciousness. It can be morally clarifying.

But the altered state is also vulnerable to inflation. When ordinary perception dissolves, the first new pattern that appears can feel like God, and sometimes it is merely your resentment wearing luminous robes. The bourgeois spa version of transformation takes the dissolution of ego as a luxury experience while leaving the structure of consumption intact. It says: I lost my self in the jungle, discovered interbeing, returned with a necklace, and increased my consulting rates. That is not liberation. That is symbolic tourism.

The deeper question is whether altered consciousness can be integrated into disciplined responsibility. Without integration, the experience becomes either entertainment or ideology. If a psychedelic experience reveals the interconnectedness of all beings but does not make you more truthful, patient, courageous, and useful to the people who actually know you, then the experience has failed, or rather you have failed the experience. The same applies to class consciousness. If it reveals exploitation but transforms you primarily into a more articulate hater, it has also failed. Hatred can diagnose injury, but it cannot build a just order.

A mature political consciousness must unite compassion for the oppressed, suspicion of exploitative structures, respect for competence, awareness of human malevolence, and caution about the tyrannical potential of utopian fantasy. That is a very difficult synthesis. It is far easier to take mushrooms, denounce capitalism, and forget to call your mother. It is also far easier to meditate yourself into a state where the suffering of others becomes "energy" you can aesthetically appreciate. Both are evasions.

So, yes, class consciousness can be an altered state if by that we mean a reorganization of perception around relations that were previously concealed. But it must be tested by action, humility, and responsibility. Otherwise it becomes a spa for the ego, even if the spa has revolutionary posters on the wall.
\end{answer}

\begin{question}
I find this useful, though I would add that hatred is not simply pathological. Sometimes hatred is the only honest form love takes when the world has made love impossible. But let us keep the tension alive. You often criticize postmodern neo-Marxism, a phrase that has become almost a floating signifier, sometimes precise, sometimes comic, sometimes like a suitcase into which all academic sins are packed. What, in your strongest formulation, is the danger? Is it that postmodernism destroys truth, that Marxism moralizes resentment, or that their hybrid produces a subject who experiences every limit as violence and every standard as domination?
\end{question}

\begin{answer}
The danger is exactly the hybrid pathology, although each component has its own dangers and its own legitimate insights. Postmodernism, at its best, tells us that categories are not innocent, that language frames perception, that power hides inside neutrality, and that the marginal voice may perceive what the center cannot. Those are important lessons. Marxism, at its best, tells us that economic structures shape consciousness, that exploitation can be systematic rather than merely personal, and that dignity requires material conditions. Those are also important lessons.

The catastrophe occurs when the suspicious hermeneutics of postmodernism fuses with the moral certainty of vulgar Marxism. Then you get an interpretive machine that can dissolve any claim to truth as a power move while simultaneously asserting its own political conclusions as morally unquestionable. That is not thought. That is acid with a gavel.

If truth is only power, then the question becomes: whose power? And if the world is fundamentally divided into oppressor and oppressed identities, then dialogue becomes confession, education becomes accusation, and disagreement becomes harm. The individual is reduced to a node in a guilt matrix. Competence becomes suspect if it is unevenly distributed. Tradition becomes nothing but domination. Beauty becomes exclusion. Standards become violence. Even grammar starts sweating under interrogation.

Now, people will say this is a caricature, and sometimes it is. There are careful scholars working on these questions with nuance. But there is also a simplified activist metaphysics that has migrated into institutions, and it is psychologically attractive because it offers moral elevation without personal transformation. You do not have to become wise; you only have to identify the contaminating structure. You do not have to repent; you only have to accuse. You do not have to forgive; forgiveness becomes complicity.

The deepest problem is spiritual. The hybrid ideology externalizes evil. It locates malevolence primarily in systems, groups, discourses, histories, and identities, and while evil certainly operates there, the doctrine becomes corrupt when it refuses to see evil in the self. The biblical warning is that the line between good and evil runs through the heart. The ideological warning label says: consult your sociological location for moral status. That is a downgrade.

But I will say this, because it matters: the critique of this pathology must not become an excuse for complacency. Real injustice exists. Cruelty hides in tradition. Exploitation hides in markets. Exclusion hides in standards. Hypocrisy hides in moral language. If conservatives refuse to see that, they deserve the postmodern critique. If radicals refuse to see the necessity of truth, competence, gratitude, and restraint, they deserve the ruins they keep accidentally designing. The task is to rescue the true fragments from the ideological blender.
\end{answer}

\begin{question}
Here maybe we agree too much, which is suspicious. Agreement is often where the unconscious puts the knife. Let me ask about pure ideology. In my work I have often insisted that ideology is not simply a false belief. It is embodied in practices, rituals, jokes, food, architecture, bureaucracy, even the little pleasures we take in knowing very well what we are doing and doing it anyway. In this sense, perhaps "clean your room" is also pure ideology, but not therefore false. It works because it organizes enjoyment. Can an ethic succeed without organizing enjoyment? Or does every moral order need a libidinal economy, some pleasure in obedience, some obscene supplement?
\end{question}

\begin{answer}
Every moral order has a libidinal economy. Of course it does. Human beings are not disembodied propositions. They are creatures of hunger, shame, pride, lust, play, rhythm, imitation, disgust, awe, and the strange pleasure of finally putting a thing where it belongs. If your ethic cannot organize enjoyment, it will be defeated by the first carnival that comes along with drums.

This is why puritanical systems fail and libertine systems fail in opposite directions. The puritan tries to build order by crushing pleasure. The result is resentment, hypocrisy, and eventually an eruption of the repressed that arrives wearing either a secret costume or a political uniform. The libertine tries to build freedom by enthroning pleasure. The result is dissociation, boredom, addiction, and the collapse of meaning into stimulation. A viable moral order does not abolish enjoyment. It educates it.

Ritual is educated enjoyment. Dance, feast, fasting, liturgy, sport, courtship, scholarship, craft, even proper argument, these are ways of binding affect to form. They allow the body to participate in order without reducing order to mere coercion. A child does not learn civilization by reading Kant. A child learns civilization by being held, fed, corrected, sung to, praised, constrained, laughed with, and gradually invited into games where rules make joy possible. That is culture at the developmental level.

"Clean your room" organizes enjoyment because order is not merely duty. There is pleasure in seeing the floor. There is relief in finding the book. There is dignity in making the bed badly and then slightly better. The danger, as you say, is that this pleasure can become ideological. One can enjoy order sadistically. One can turn tidiness into contempt for those in chaos. One can polish the surface while the heart becomes a basement full of rats. But the danger of corruption does not erase the moral fact that disciplined pleasure is necessary.

The obscene supplement is unavoidable because every moral code produces excess. The stated rule says: be responsible. The supplement whispers: and feel superior to the irresponsible. The stated rule says: tell the truth. The supplement whispers: use truth to humiliate the liar. The stated rule says: care for the poor. The supplement whispers: display your caring as proof of your innocence. That is why moral systems require confession, comedy, friendship, and enemies who are allowed to speak. Without mechanisms for exposing the obscene supplement, virtue becomes narcissism with better posture.

So yes, ethics must organize enjoyment. But it must also humble enjoyment. The mature form of pleasure is not immediate gratification but participation in a pattern that can survive contact with truth. That is why music is such a profound model. Constraint produces freedom. Rhythm permits improvisation. The orchestra is hierarchical, cooperative, disciplined, ecstatic, and prone to catastrophe if the percussion section decides that standards are oppressive. That is civilization in miniature.
\end{answer}

\begin{question}
Since we are now in music, let me free jazz the question in the vulgar way: what if consciousness itself is not the melody but the mistake that makes the melody audible? In Hegelian terms, spirit appears through negativity; in Lacanian terms, the subject appears where the signifier fails; in Marxian terms, class appears where the smooth reproduction of society breaks down. In neuroscience, perhaps error signals do the same thing. So perhaps consciousness is not order against chaos, but the experience of a glitch in order. Does your emphasis on order sufficiently honor the productive dignity of breakdown?
\end{question}

\begin{answer}
Breakdown is often the birthplace of consciousness. I agree. You do not become conscious of the hammer until it misses the nail and hits your thumb. You do not become conscious of the body until it hurts. You do not become conscious of trust until betrayal. You do not become conscious of speech until you cannot say what you mean, which, for some people, becomes an academic career.

Order that never breaks becomes unconscious habit. That is not bad. Habit is a miracle. If you had to consciously calculate how to walk, speak, breathe, and interpret facial expression every moment, you would collapse under the computational burden before breakfast. Much of sanity is benevolent unconscious order. But when order fails, attention awakens. Error is the summons. That is the predictive-processing insight in moral clothing: the world violates expectation, and the organism must update or perish.

The productive dignity of breakdown depends on whether the breakdown is integrated. Some breakdowns destroy. Trauma can reveal reality, but it can also fragment the person so severely that revelation becomes unusable. Psychosis may disclose the symbolic instability of ordinary life, but it may also devour the sufferer. Political crisis may reveal structural injustice, but it may also unleash the worst aspects of group psychology. So we must not romanticize breakdown. The abyss is not a graduate seminar with better lighting.

But neither should we pathologize all breakdown. Depression can be a signal that a mode of life has become untenable. Anxiety can indicate that the future is demanding preparation. Anger can reveal violated boundaries. Grief can enlarge the soul by forcing it to love what it can no longer possess. Doubt can rescue faith from stupidity. These are not pleasant processes. They are not "wellness." They are descents.

My emphasis on order is precisely because breakdown is powerful. If you are going to enter chaos, you need a vessel. Ritual, discipline, friendship, skill, language, tradition, and voluntary responsibility provide the container in which breakdown can become transformation rather than mere disintegration. The alchemists understood this symbolically. You dissolve and coagulate. You do not merely dissolve. Pure dissolution is acid. Pure coagulation is stone. Consciousness matures by cycling between them.

So, yes, consciousness is in part the glitch that makes the melody audible. But if there is no melody, the glitch is only noise. And if there is no capacity to repair, the glitch becomes identity. That is another modern danger: people fall in love with the site of injury because it gives them a place to speak from. But the deepest speech is not "I am broken." It is "I was broken there, and from that fracture I learned how form can be merciful." That is not motivational optimism. It is tragic realism with a broom.
\end{answer}

\begin{question}
We should speak about the father, because otherwise the father will speak through us anyway. In your work, the father often appears as discipline, logos, tradition, law, the voice that says: stand up straight, tell the truth, enter the world. In psychoanalysis, of course, the Name-of-the-Father is not simply daddy with a rulebook. It is the symbolic function that cuts the child from immediate enjoyment. But what happens today when the paternal function is outsourced to institutions, algorithms, therapy culture, online audiences, and market metrics? Who is the father in a world where the superego says not "obey" but "optimize"?
\end{question}

\begin{answer}
The father becomes dispersed, and therefore more difficult to confront. That is one of the central pathologies of our moment. In a traditional order, for all its faults, the law had a face. The father, the priest, the elder, the teacher, the judge, the craftsman, the story, the rite of passage, these were imperfect but embodied mediators of constraint. You could rebel against them, learn from them, forgive them, or become them. Today the law is often faceless. It appears as metrics, screens, policies, feeds, ratings, institutional procedures, algorithmic nudges, and therapeutic imperatives. It does not say "obey" in the old voice. It says "curate yourself."

The optimizing superego is far more tyrannical than the forbidding father because it has no sabbath. The old law said: do not. The new law says: become endlessly better at being yourself in a way that is measurable, visible, monetizable, emotionally transparent, politically acceptable, physically attractive, economically productive, and available for feedback. That is not liberation. That is a treadmill with a mirror.

The paternal function, properly understood, is not domination. It is the merciful introduction of limit. A good father says: you cannot be everything, and this is good news, because now you can become something. He says: your impulses are not the same as your destiny. He says: the world will resist you, and I will help you become strong enough that its resistance does not turn you bitter. He says: no, and the no creates the boundary in which a yes can grow.

When that function is outsourced to algorithms, the limit becomes inhuman. The algorithm does not love you. It does not want your maturity. It wants engagement, and engagement is often generated by envy, outrage, lust, terror, and humiliation. A culture raised by algorithms becomes fathered by appetite. Therapy culture can help, but it can also become maternalized without limit, endlessly validating the self's narrative without demanding transformation. Institutions can help, but they increasingly speak in procedural language that conceals moral cowardice. Markets can coordinate value, but they cannot tell you what is sacred.

So who is the father now? Potentially, no one, which means everyone becomes haunted by a demand they cannot locate. Or, more hopefully, the paternal function must be consciously regenerated in families, schools, friendships, crafts, religious communities, and serious mentorship. This does not mean returning to authoritarian caricature. It means recovering loving constraint. Without it, young people are left alone with infinite possibility, which is indistinguishable from chaos when no one has taught you how to choose.
\end{answer}

\begin{question}
Your "loving constraint" brings us to religion. For some, religion is the symbolic immune system of civilization; for others, the most successful ideological hallucination in human history. I want to avoid the stupid debate between belief and disbelief. The more interesting question is whether consciousness requires sacred fiction. Hegel would say religion is spirit representing itself to itself in picture-thinking. Lacan would say God is located in the structure of the big Other, even when we say the big Other does not exist. Marx would mutter something about opium, but opium also relieves pain. So is religion an adaptive myth, a metaphysical truth, a neurocognitive prosthesis, or a necessary lie that became too deep to call a lie?
\end{question}

\begin{answer}
Religion is what happens when a culture encodes the conditions for psychological and social survival in stories deep enough to outlast the people who first understood them. That is one answer. It is not the whole answer, but it is a necessary beginning.

If you reduce religion to false propositions about the physical world, you miss almost everything. If you reduce religion to metaphysical certainty, you also risk missing the existential function. Religious stories are not scientific hypotheses in primitive clothing. They are symbolic maps of value. They tell you what to aim at, how to suffer, how to repent, how to forgive, how to marry, how to bury your dead, how to restrain power, how to celebrate abundance without becoming a pig, and how to understand the voice inside you that knows when you are lying.

Now, is that adaptive myth, metaphysical truth, prosthesis, or necessary lie? It may be all of those, except I would resist "lie" because a lie is a deliberate distortion of truth for manipulative purposes. A myth is not a lie. A myth is a pattern of meaning represented in narrative form. The fact that it is not literal in the narrow empirical sense does not make it false. If a child says the story of the dragon taught him courage, and a professor replies that dragons are taxonomically under-specified, the professor is the fool.

Consciousness probably does require sacred fiction in the sense that finite beings need representations of ultimate value that exceed immediate utility. But if the sacred fiction is merely invented, merely arbitrary, it will not sustain sacrifice. People can smell fabrication, even when they pretend not to. The sacred must be experienced as discovered, not merely constructed. This is where the modern mind struggles. It wants the benefits of transcendence without the humility of receiving anything. It wants to author the sacred and then be transformed by it, which is like writing yourself a love letter and being surprised by the handwriting.

As for metaphysical truth, I would say the religious impulse points toward realities that are not exhausted by material description. Whether one cashes that out theologically, phenomenologically, or pragmatically is a profound question. But the attempt to live without the sacred does not produce neutrality. It produces substitute religions: politics, nation, race, sex, therapy, markets, technology, self-expression, health, apocalypse, even anti-religion itself. The sacred does not disappear. It migrates, often into cruder vessels.

So religion is not simply opium. Sometimes it is morphine for real agony. Sometimes it is poison. Sometimes it is discipline. Sometimes it is the memory of civilization encoded as song. Sometimes it is tyranny with incense. But at its best, religion is the dramatization of the proposition that truth, love, and voluntary sacrifice are more real than appetite, power, and death. That proposition may be the deepest antientropic structure consciousness has yet discovered.
\end{answer}

\begin{question}
We are approaching the end, which means, dialectically, that we must return to the beginning but with more confusion. The lobster, the room, the father, entropy, class consciousness, postmodernism, religion: perhaps all these are different names for the same impossible mediation between the body and the symbolic. Let me ask in the most ALIUS way: what altered state of consciousness would allow humanity to experience hierarchy without domination, equality without ressentiment, truth without cruelty, myth without fascism, and entropy without despair? Or is the desire for such a state already the fantasy of a final reconciliation that would abolish consciousness itself?
\end{question}

\begin{answer}
The desire for final reconciliation is both necessary and dangerous. It is necessary because without an image of reconciliation, however impossible, we cannot orient ourselves toward peace. It is dangerous because if we believe we can instantiate final reconciliation politically, psychologically, or technologically, we will start sacrificing actual human beings to the imagined harmony of the system. Paradise becomes most lethal when it is given an implementation schedule.

The altered state you describe would not be ecstasy in the usual sense. It would be something like lucid humility. That is not a phrase people will put on a nightclub poster, but perhaps they should. Lucid humility would be the state in which one perceives hierarchies as inevitable but morally accountable, equality as sacred but not flattening, truth as liberating but wounding, myth as necessary but corruptible, entropy as tragic but meaning-giving, and the self as responsible but not sovereign. It would be a state in which complexity does not produce cynicism and simplicity does not become stupidity.

How does one cultivate such a state? Not primarily by pharmacology, although altered states may help some people glimpse dimensions of experience normally hidden. Not primarily by ideology, because ideology tends to harden the glimpse into a weapon. Not primarily by therapy, because therapy can become a circular polishing of the self. The cultivation requires practices: truthful speech, disciplined attention, embodied service, ritual, study, art, friendship, courageous conflict, forgiveness, competence, and the willingness to be corrected by reality before reality becomes catastrophic.

Hierarchy without domination requires that rank be tied to service and competence rather than mere power. Equality without ressentiment requires gratitude for the being of others and a refusal to make excellence into an insult. Truth without cruelty requires love, because truth without love becomes a scalpel in the hands of a sadist, while love without truth becomes anesthesia. Myth without fascism requires pluralism, humor, and the constant remembrance that the hero's shadow is the tyrant. Entropy without despair requires that we learn to see finitude not merely as defeat but as the condition that gives urgency to care.

Will humanity achieve this? Not finally. Final achievement would abolish the drama. Consciousness is not meant to end in a perfectly reconciled spreadsheet of the soul. It is meant to participate in the ongoing transformation of chaos into habitable order, knowing that the transformation will never be complete. That incompletion is not a defect. It is the opening through which freedom enters.

So we return to the lobster. It climbs because its body tells it that ascent matters. We speak because our suffering tells us that ascent must be justified. Between the claw and the word lies the whole tragedy of consciousness. And perhaps the proper altered state is not to escape that tragedy, but to inhabit it with enough courage that our hierarchies become less cruel, our equality less resentful, our myths less murderous, our rooms less symbolic of apocalypse, and our laughter less contemptuous. That would be a good beginning. Not the end of entropy. A decent arrangement of furniture before the next storm.
\end{answer}

\begin{question}
One final, final question, which means it is not final, because finality is the fantasy by which the interviewer avoids responsibility. If a young researcher in consciousness studies reads this imaginary exchange and asks, "What should I study now?", what would you tell them? The neuronal correlates? Psychedelics? Predictive processing? Ritual? Political economy? Lobster posture? The phenomenology of cleaning? Or the strange zone where all these absurdly touch?
\end{question}

\begin{answer}
Study the strange zone where they touch, but do not do so lazily. Interdisciplinarity is often the word people use when they have not mastered one discipline but would like honorary access to six. You need rigor somewhere. Learn a method. Learn statistics, ethnography, neuroimaging, phenomenology, clinical practice, philosophy, computational modeling, textual interpretation, something. Become genuinely useful in a domain where error can find you.

Then, once you have a spine, look across the levels. Consciousness will not be understood from one level alone. The molecular matters. Receptors matter. Sleep matters. Development matters. Trauma matters. Attention matters. Culture matters. Ritual matters. Language matters. Political economy matters. Myth matters. The posture of the defeated body matters. The story that body tells about its defeat matters. The institution that profits from the defeat matters. The drug that briefly alters the salience of the defeat matters. The friend who says "stand up" matters. The fact that "up" is already a metaphysical direction matters.

But beware of total theories. A total theory of consciousness is very tempting because consciousness is what totality feels like from the inside. The temptation will be to say: it is all prediction, or it is all power, or it is all serotonin, or it is all language, or it is all free energy, or it is all trauma, or it is all archetype. Each "all" is a little idol. Useful at first, then demanding sacrifices.

The young researcher should cultivate methodological pluralism and existential seriousness. Ask clear questions. Do not use obscurity to protect weak claims, unless you are writing satire, in which case make the obscurity so excessive that it becomes diagnostically useful. Respect data, but do not worship it. Respect subjective reports, but do not drown in them. Respect biology, but do not use it to excuse cruelty. Respect social critique, but do not use it to flee your own conscience.

And finally, study what frightens you intellectually. If hierarchy frightens you, study it without reducing it to oppression. If equality frightens you, study it without reducing it to envy. If religion frightens you, study it without reducing it to superstition. If materialism frightens you, study it without reducing it to nihilism. If postmodernism frightens you, read the best version rather than the campus slogan. If order frightens you, clean one shelf. If chaos frightens you, read a poem you do not understand.

The strange zone is where consciousness lives because consciousness is precisely the mediator among levels that do not know how to speak to each other. It is lobster serotonin dreaming of Hegel. It is a Marxist analysis with a nervous system. It is a ritual trying to become a receptor model. It is a room becoming a cosmos by way of laundry. Study that, and keep your claims modest enough that reality can improve them.
\end{answer}

\section*{Selected Excerpts}
\begin{itemize}
\item Consciousness is not simply a ladder. It is the capacity to suffer the ladder, to resent it, to climb it, to discover that it was partly imaginary, and then to build a better one.
\item The end of entropy would be hell disguised as accounting.
\item The TED talk is the archetype after it has been bitten by the commodity form and then forced to smile under stage lighting for eighteen minutes.
\item The optimizing superego is far more tyrannical than the forbidding father because it has no sabbath.
\item Between the claw and the word lies the whole tragedy of consciousness.
\end{itemize}

\section*{References}
\begingroup
\small
Baars, B. J. (1988). \textit{A cognitive theory of consciousness}. Cambridge University Press.\par
Carhart-Harris, R. L., and Friston, K. J. (2019). REBUS and the anarchic brain: Toward a unified model of the brain action of psychedelics. \textit{Pharmacological Reviews, 71}(3), 316-344.\par
Dehaene, S. (2014). \textit{Consciousness and the brain: Deciphering how the brain codes our thoughts}. Viking.\par
Friston, K. (2010). The free-energy principle: A unified brain theory? \textit{Nature Reviews Neuroscience, 11}(2), 127-138.\par
Hegel, G. W. F. (1977). \textit{Phenomenology of spirit} (A. V. Miller, Trans.). Oxford University Press. (Original work published 1807)\par
Lacan, J. (1977). \textit{Ecrits: A selection} (A. Sheridan, Trans.). Norton.\par
Marx, K. (1976). \textit{Capital: A critique of political economy, Volume 1} (B. Fowkes, Trans.). Penguin. (Original work published 1867)\par
Peterson, J. B. (1999). \textit{Maps of meaning: The architecture of belief}. Routledge.\par
Peterson, J. B. (2018). \textit{12 rules for life: An antidote to chaos}. Random House Canada.\par
Seth, A. K. (2021). \textit{Being you: A new science of consciousness}. Faber and Faber.\par
Zizek, S. (1989). \textit{The sublime object of ideology}. Verso.\par
Zizek, S. (2006). \textit{The parallax view}. MIT Press.\par
\endgroup

\end{document}
